\section{Mejoras a futuro}
Para finalizar el anteproyecto se presentan esta última sección que mencionan todas aquellas cosas que se quieren realizar pero que por un motivo u otro quedan fuera del alcance.

\subsection{Diseño de GUI para usuario}
Una GUI es una interfaz gráfica que permite al usuario poder manipular el dispositivo conectado desde la PC. Esta misma haría uso del conector Ethernet que se encuentra dentro del propio dispositivo, el mismo que permite el uso para laboratorio remoto. La idea detrás de este GUI sería lograr que el usuario pueda realizar las mismas configuraciones que venía haciendo para la generación de la forma de onda, pero además permita cargar tablas de muestras para una nueva forma arbitraria definida por el propio usuario.

\subsection{Implementación de funcionalidades/contador}
A modo de comparar con el equipo proporcionado por FeelElec FY6900 queda como futura mejora la implementación de funcionalidades extras que permitan darle un mayor valor al propio instrumento realizado, así como también la implementación de un contador digital de frecuencia que puede realizarse con una FPGA.

\subsection{Mejora del display con LVGL}
La idea de esta mejora consiste en la utilización de un nuevo display (ya contemplado en el diseño del hardware) que permita el uso de librerías como LVGL, de forma que brinde al usuario una mejor representación gráfica.

\subsection{Diseño de una carcasa y frente adecuados}
Debido a la longitud del proyecto algunas cosas como la presentación final del instrumento no sean adecuadas, la idea de esta mejora consiste en realizar una carcasa y frente a nivel industrial, de forma que el instrumento pueda ser comerciable y adquirido por cualquier usuario.
